
\section{Introduction}
\label{sec:intro}

Industrial workers in manufacturing, foundry, and chemical-processing sectors routinely
endure elevated exposure to respirable particulate matter (\pmone, \pmtf, \pmten), carbon
monoxide (CO), nitrogen oxides (NO$_x$), hydrogen sulfide (H$_2$S), and volatile organic
compounds (VOCs)~\citep{nasri2023pm25,sekhavati2023particulate}. Chronic inhalation of
fine particulates below 2.5~$\mu$m aerodynamic diameter is associated with alveolar
inflammation, cardiovascular disease, pulmonary fibrosis, and impaired cognitive
health~\citep{pai2022who}. In Bangladesh, industrial corridors such as Tejgaon, Gazipur,
and Narayanganj routinely register ambient \pmtf{} concentrations exceeding 150~\ugsm{}
during winter inversions, and within active factory microenvironments pollutant levels can
increase substantially above background within seconds of process
initiation~\citep{basak2025spatial,ali2025longterm}.

Conventional centralized ventilation requires prohibitive capital expenditure and
multi-kilowatt continuous operation, making it economically unviable for the majority of
SMEs in developing economies~\citep{chowdhury2025misbeliefs}. Low-cost IoT sensing
architectures have emerged as practical tools for fine-grained environmental
surveillance~\citep{alam2025iot}; however, two recognized deficiencies persist.
First, existing commercial purifiers activate \textit{only after} particulate concentrations
breach hazardous thresholds---by which time the transition from clean air to hazardous
exposure has already occurred~\citep{dubey2021assessing}. Second, static batch-trained
forecasting models can suffer performance degradation when relocated across disparate
factory zones due to concept drift~\citep{noman2025aiair}, and lack on-site adaptation
capacity without batch retraining. Recent work has begun exploring predictive
purifier control~\citep{beldar2025aipurifier,energyeff2024benv} and incremental learning
for air quality monitoring~\citep{incremental2024asc}, though closed-loop integration
of online learning with localized physical purification in ultra-low-cost hardware
remains incompletely addressed.

\begin{figure}[htbp]
\centering
\includegraphics[width=0.93\linewidth]{system_workflow.png}
\caption{NirmalNode closed-loop workflow: edge sensing $\to$ EMA filtering $\to$ 14-feature
  vector $\to$ Adaptive AI Model Arena $\to$ Predictive Decision Engine $\to$ GPIO relay
  $\to$ three-stage purification, with incremental model update closing the loop.}
\label{fig:arch}
\end{figure}

This paper presents \textbf{NirmalNode}, a closed-loop Green IoT architecture coupling
multi-sensor hotspot monitoring, online incremental stream forecasting, predictive
relay actuation, and localized three-stage physical filtration in a single autonomous
node (Fig.~\ref{fig:arch}). The core contributions are:
\begin{itemize}[leftmargin=*,topsep=2pt,itemsep=1pt]
  \item \textbf{Integrated closed-loop edge--cloud node:} an autonomous IoT node coupling
    multi-channel aerosol/gas sensing with a three-stage physical filtration train and
    a direct connection between forward prediction and electromechanical actuation within
    an ultra-low-cost (\$67.27 prototype BOM) footprint.
  \item \textbf{Online stream-learning Model Arena:} nine parallel incremental estimators
    under a prequential (test-then-train) protocol, electing an adaptive SGD+PA ensemble
    that updates online to site-specific observations without manual model retraining.
  \item \textbf{Predictive purifier control engine:} proactive relay actuation based on
    short-horizon PM$_{2.5}$ forecasts, providing measurable actuation lead time before
    the monitored concentration crosses the selected threshold.
  \item \textbf{Multi-stream empirical evaluation:} zero missed hazard events on a
    Bangladesh public-dataset-derived evaluation stream and 95.7\% single-pass \pmtf{}
    reduction in a laboratory aerosol bench trial.
\end{itemize}
